\scanpage{436}
\mspage{3}
% archive sheet number: 435

Ainsi on trouve bien un foncteur (en fait, \unsure{$=$|même} un 2-foncteur de 2-catégories)

\[
f_!\colon\uline{\mathrm{Grp}}(S)\longrightarrow\uline{\mathrm{Grp}}(T)
\]

Si

\[
u\colon G\longrightarrow G'
\]

est un \unsure{hom|lien} de groupoïdes fibrés sur $S$, qui induise \deleted{des} \inserted{Cat-}équivalences sur les fibres, alors

\[
f_!(u)\colon f_!(\unsure{G|u})\unsure{\mathrel{-}|\longrightarrow}f_!(G')
\]

est aussi une Cat-équiv sur les fibres. (Cela est formel, à partir du fait que $f_!$ est un 2-foncteur, compte tenu que si on a une équiv. $u$ sur les fibres, alors il y a un quasi-inverse -- c'est une histoire générale de cat. fibrées\ldots)

Soit maintenant

\[
p\colon X\longrightarrow S,
\]

et prenons le cas $G=\Pi(X/S)$. Je dis qu'on a alors une \uline{$T$-équivalence} \unsure{can}

\[
\boxed{f_!\bigl(\Pi(X/S)\bigr)\simeq\deleted{\unsure{g}}\Pi(X/T)}
\]

en d'autres termes, pour connaître $\Pi(X,T)$ \inserted{(à $T$-équiv près)}, il suffit de connaître $\Pi(X/S)$, et de lui appliquer le foncteur $f_!$ d'intégration des 1-types d'homotopie relatifs, pour $f\colon S\longrightarrow T$.

Notons que l'on a trivialement

\[
\Pi(X/S)=p_!(\unsure{e}_X)
\]

où $e_X$ désigne le $X$-groupoïde fibré trivial (toutes les fibres ponctuelles). Donc la formule précédente s'écrit comme une formule de transitivité

\[
f_{\deleted{\unsure{*}}!}\bigl(p_!(\unsure{\Pi|\Gamma|\mathscr T})\bigr)\simeq(fp)_!(\unsure{\Gamma|\mathscr T|\mathrm{I}}),
\]
